\documentclass[12pt]{article}
\usepackage{amsmath, amssymb, amsthm}
\usepackage{geometry}
\geometry{margin=1in}

\title{CSE 400: Fundamentals of Probability in Computing\\
Tutorial-1 Lecture Scribe}
\date{February 13, 2026}

\begin{document}
\maketitle

\section*{Q1: Partitioning of Dishes}

\subsection*{Problem}
20 distinct dishes are divided into 4 unlabeled groups of 5 dishes each.

\subsection*{Result (a)}
The number of ways to divide the dishes is
\[
\frac{20!}{(5!)^4 4!}.
\]

\subsection*{Derivation}
Since:
\begin{itemize}
    \item All 20 dishes are distinct,
    \item Each group contains 5 dishes,
    \item Order within each group does not matter,
    \item Order of groups does not matter,
\end{itemize}
we apply the multinomial coefficient formula.

\subsection*{Result (b)}
If each cuisine contains exactly 5 dishes, the probability that each platter contains dishes from only one cuisine is:
\[
\frac{(5!)^4 4!}{20!}.
\]

\section*{Q2: Uniform Arrival and Waiting Time}

\subsection*{Assumption}
Arrival time is uniformly distributed over 30 minutes.

\subsection*{(a) Waiting time < 5 minutes}
Favorable intervals: 10 minutes.

\[
P(\text{wait} < 5) = \frac{10}{30} = \frac{1}{3}.
\]

\subsection*{(b) Waiting time > 10 minutes}
Favorable intervals: 10 minutes.

\[
P(\text{wait} > 10) = \frac{10}{30} = \frac{1}{3}.
\]

\section*{Q3: Conditional Probability}

\subsection*{Given}
\[
P(L)=0.15,\quad P(E)=0.25,\quad P(L\cap E)=0.08.
\]

\subsection*{Compute}
\[
P(E^c)=0.75.
\]

\[
P(L\cup E)=0.15+0.25-0.08=0.32.
\]

\[
P(L^c \cap E^c)=1-0.32=0.68.
\]

\[
P(L^c|E^c)=\frac{0.68}{0.75}\approx 0.907.
\]

\section*{Q4: Poisson Distribution (λ=0.2)}

\subsection*{Definition}
\[
P(X=k)=e^{-\lambda}\frac{\lambda^k}{k!}.
\]

\subsection*{(a) P(X=0)}
\[
P(X=0)=e^{-0.2}=0.8187.
\]

\subsection*{(b) P(X\ge2)}
\[
P(X\ge2)=1-P(0)-P(1).
\]

\[
P(1)=0.2e^{-0.2}=0.1637.
\]

\[
P(X\ge2)=1-0.8187-0.1637=0.0176.
\]

\section*{Q5: Poisson Distribution (λ=3.5)}

\subsection*{(a) P(X\ge2)}
\[
P(X\ge2)=1-P(0)-P(1).
\]

\[
P(0)=0.0302,\quad P(1)=0.1057.
\]

\[
P(X\ge2)=0.8641.
\]

\subsection*{(b) P(X\le1)}
\[
P(X\le1)=0.1359.
\]

\section*{Q6: Poisson Distribution (λ=3)}

\subsection*{(a) P(X\ge3)}
\[
P(X\ge3)=1-P(0)-P(1)-P(2)=0.5768.
\]

\subsection*{(b)}
\[
P(X\ge3|X\ge1)=\frac{0.5768}{0.9502}\approx 0.607.
\]

\section*{Q7: Gaussian Distribution}

\subsection*{PDF}
\[
f_X(x)=\frac{1}{\sqrt{8\pi}}\exp\left(-\frac{(x+3)^2}{8}\right).
\]

\[
m=-3,\quad \sigma=2.
\]

\subsection*{Results}
\[
P(X\le0)=\Phi\left(\frac{3}{2}\right)=1-Q\left(\frac{3}{2}\right).
\]

\[
P(X>4)=Q\left(\frac{7}{2}\right).
\]

\[
P(|X+3|<2)=1-2Q(1).
\]

\[
P(|X-2|>1)=1-Q(2)+Q(3).
\]

\section*{Q8: Jury Decision}

If defendant is guilty:
\[
\sum_{i=8}^{12} \binom{12}{i}\theta^i(1-\theta)^{12-i}.
\]

If innocent:
\[
\sum_{i=5}^{12} \binom{12}{i}\theta^i(1-\theta)^{12-i}.
\]

Total probability:
\[
\alpha \sum_{i=8}^{12} \binom{12}{i}\theta^i(1-\theta)^{12-i}
+
(1-\alpha)\sum_{i=5}^{12} \binom{12}{i}\theta^i(1-\theta)^{12-i}.
\]

\section*{Q9: PMF Normalization}

\[
p(i)=c\frac{\lambda^i}{i!}.
\]

\[
ce^\lambda=1 \Rightarrow c=e^{-\lambda}.
\]

\[
P(X=0)=e^{-\lambda}.
\]

\[
P(X>2)=1-e^{-\lambda}-\lambda e^{-\lambda}-\frac{\lambda^2}{2}e^{-\lambda}.
\]

\section*{CDF Definition}

\[
F(a)=\sum_{x\le a} p(x).
\]

Example:
\[
F(a)=
\begin{cases}
0, & a<1 \\
\frac14, & 1\le a<2 \\
\frac34, & 2\le a<3 \\
\frac78, & 3\le a<4 \\
1, & a\ge4
\end{cases}
\]

\section*{Q10: System Reliability}

\subsection*{(a)}
5-component system better than 3-component system if:
\[
3(p-1)^2(2p-1)>0
\Rightarrow p>\frac12.
\]

\subsection*{(b)}
\[
P_{2k+1}-P_{2k-1}
=
\binom{2k-1}{k} p^k(1-p)^k(2p-1).
\]

\[
>0 \iff p>\frac12.
\]

\end{document}